\documentclass[a4paper,12pt]{article}
\usepackage[T2A]{fontenc}
\usepackage[utf8]{inputenc}
\usepackage[russian]{babel}
\usepackage{amsmath,amssymb,booktabs,tabularx,array,xcolor,tikz,fancyhdr,geometry}
\usepackage[colorlinks=true,linkcolor=teal]{hyperref}
\geometry{left=20mm,right=20mm,top=23mm,bottom=22mm,headheight=15pt}
\definecolor{accent}{HTML}{006D70}
\setlength{\parindent}{0pt}
\setlength{\parskip}{6pt}
\renewcommand{\arraystretch}{1.35}
\pagestyle{fancy}
\fancyhf{}
\fancyhead[L]{\small Ревью домашней работы}
\fancyhead[R]{\small Ярослав Верещагин}
\fancyfoot[C]{\thepage}
\renewcommand{\headrulewidth}{0.4pt}
\newcommand{\task}[2]{\clearpage\phantomsection\label{task:#1}{\Large\color{accent} Задание №~#1. #2}\par\vspace{3pt}{\color{accent}\hrule}\vspace{6pt}}
\newcommand{\parttitle}[1]{\par\textbf{#1.}\enspace}
\newcommand{\err}[1]{\par\textcolor{red!75!black}{\textbf{Ошибка:}} #1}
\newcommand{\idea}[1]{\par\textcolor{blue!60!black}{\textbf{Ключевая идея:}} #1}
\newcommand{\result}[1]{\par\textcolor{teal!60!black}{\textbf{Итог:}} #1}
\newcommand{\practice}[1]{\par\medskip{\color{accent}\textbf{Задача на закрепление}}\par #1}
\begin{document}
\begin{titlepage}
\centering
\vspace*{25mm}
{\large\color{accent} МАТЕМАТИКА\par}
\vspace{22mm}
{\Huge Ревью домашней работы\par}
\vspace{15mm}
{\Large Ярослав Верещагин\par}
\vspace{9mm}
{\large 26 ноября 2026 г.\par}
\vfill
{\large Лёвин Артём Александрович\par}
{эксклюзивно для Ярослава Верещагина\par}
\vspace{12mm}
\begin{tikzpicture}
\draw[accent!65,line width=.7pt] (0,0) .. controls (3,1.2) and (7,-1.2) .. (12,0);
\end{tikzpicture}
\end{titlepage}
\setcounter{page}{1}
{\Large\color{accent} Сводная таблица}\par\vspace{4pt}{\color{accent}\hrule}\medskip
В таблице представлены только задания, в которых необходимо исправить решение или дополнить его.

\noindent\begin{tabularx}{\linewidth}{@{}c >{\raggedright\arraybackslash}p{27mm} >{\raggedright\arraybackslash}X >{\raggedright\arraybackslash}p{27mm} >{\raggedright\arraybackslash}p{28mm}@{}}
\toprule
№ & Тема & Суть ошибки & Ваш ответ & Правильный ответ\\
\midrule
\hyperref[task:4]{4} & Квадратный корень & Некорректное ограничение корня; итог не записан & Не указан & $84$\\
\hyperref[task:7]{7} & Квадратное уравнение & Пропущен минус перед коэффициентом $b$ при записи второго корня & $-2;\ -1001$ & $-2;\ -1001$\\
\hyperref[task:10]{10} & Задача на движение & Решение на фото отсутствует & Не представлен & $18$ км/ч\\
\bottomrule
\end{tabularx}
\clearpage
\thispagestyle{empty}
\vspace*{\fill}
\begin{center}{\Huge\color{accent} Разбор ошибок}\end{center}
\vspace*{\fill}
\task{4}{Извлечение квадратного корня}
\parttitle{Текст задания} Вычислите $\sqrt{7056}$, предварительно ограничив корень соседними десятками и проверив последнюю цифру.
\parttitle{Ваше решение} Записано $80^2<\sqrt{7056}<90^2$. Далее Вы рассматриваете $84^2$ и $86^2$ и находите $84^2=7056$.
\parttitle{Ваш ответ} Не указан отдельно; найден верный квадрат $84^2$.
\err{Смешаны корень и квадраты границ.}
\parttitle{Где возникает ошибка} Вы правильно выбрали десятки 80 и 90, но уже в первом неравенстве сравнили корень с квадратами этих десятков.
\parttitle{Почему это неверно} Квадраты равны 6400 и 8100, а сам корень равен 84. Число 84 не больше 6400. При сравнении нужно либо возвести в квадрат все три неотрицательные величины, либо извлечь корень из всех трёх. Порядок неравенства при этом сохраняется.
\idea{Сначала ограничьте подкоренное число квадратами соседних десятков, затем извлеките корни из всех частей неравенства.}
\parttitle{Исправление и полное решение} Вычислим квадраты границ:
\[80^2=6400,\quad 90^2=8100,\quad 6400<7056<8100.\]
Следовательно, $80^2<7056<90^2$, откуда
\[80<\sqrt{7056}<90.\]
Квадрат оканчивается на 6 только тогда, когда целое число оканчивается на 4 или 6. Поэтому среди целых кандидатов остаются 84 и 86.
\parttitle{Проверка} Проверим первый кандидат:
\[84^2=(80+4)^2=6400+640+16=7056.\]
Число 84 положительное, поэтому это именно арифметический квадратный корень. Проверять второй кандидат уже не требуется: положительный квадратный корень единственен.
\result{$\sqrt{7056}=84$. Запишите ответ явно.}
\practice{Вычислите $\sqrt{5776}$, предварительно ограничив корень соседними десятками и проверив последнюю цифру.}
\task{7}{Квадратное уравнение}
\parttitle{Текст задания} Решите уравнение, используя удобный способ извлечения корня из дискриминанта: $x^2+1003x+2002=0$.
\parttitle{Ваше решение} Найден дискриминант $D=998001$ и его корень 999. Первый корень записан как $(-1003+999)/2=-2$. В числителе второго корня перед 1003 не виден знак минус.
\parttitle{Ваш ответ} $-2;\ -1001$.
\err{Ответ верный, но запись второго корня не согласуется с формулой.}
\parttitle{Где возникает ошибка} Последний верный шаг --- вычисление первого корня. При записи второго корня необходимо сохранить $-1003$; меняется только знак перед корнем из дискриминанта.
\parttitle{Почему это неверно} Для $ax^2+bx+c=0$ при $a\ne0$ и $D>0$ имеем
\[x_1=\frac{-b+\sqrt D}{2a},\qquad x_2=\frac{-b-\sqrt D}{2a}.\]
Знак перед $b$ одинаков в обеих формулах. Если вместо $-1003$ написать $1003$, то $(1003-999)/2=2$, а не $-1001$.
\idea{В обоих числителях сначала пишите $-b$, затем выбирайте знак перед $\sqrt D$.}
\parttitle{Исправление и полное решение} Здесь $a=1$, $b=1003$, $c=2002$.
\[D=1003^2-4\cdot2002=1006009-8008=998001.\]
Заметим, что
\[999^2=(1000-1)^2=1000000-2000+1=998001.\]
Значит, $\sqrt D=999$. Получаем
\[x_1=\frac{-1003+999}{2}=-2,\qquad x_2=\frac{-1003-999}{2}=-1001.\]
\parttitle{Проверка} По теореме Виета сумма корней равна $-2-1001=-1003$, произведение равно $(-2)(-1001)=2002$. Оба условия выполнены.
\result{Ответ $-2;\ -1001$ верен. Исправьте знак во втором числителе.}
\practice{Решите $x^2+503x+1002=0$, используя удобный способ извлечения корня из дискриминанта.}
\task{10}{Задача на движение}
\parttitle{Текст задания} Между пунктами $A$ и $B$ --- 36 км. Обратно автомобиль едет на 6 км/ч быстрее, затрачивая на 1 ч меньше. Найдите скорость на обратном пути. Составьте таблицу и объясните, почему отрицательный корень не подходит.
\parttitle{Ваше решение и ответ} На фотографиях отсутствуют. Поэтому нельзя установить, какой способ Вы использовали и где могли ошибиться.
\err{Не представлено решение; обязательные таблица и обоснование выбора корня отсутствуют на фото.}
\idea{Выразите обе скорости через одну переменную, найдите время по формуле $t=S/v$ и составьте уравнение по разности времён.}
\parttitle{Полное решение} Пусть $x$ км/ч --- скорость на обратном пути. Тогда скорость туда равна $x-6$ км/ч. Обе скорости положительны, поэтому $x>6$.
\begin{center}
\begin{tabular}{lccc}
\toprule
Путь & $S$, км & $v$, км/ч & $t$, ч\\
\midrule
Туда & 36 & $x-6$ & $\dfrac{36}{x-6}$\\
Обратно & 36 & $x$ & $\dfrac{36}{x}$\\
\bottomrule
\end{tabular}
\end{center}
Время поездки туда на один час больше времени обратной поездки. Следовательно,
\[\frac{36}{x-6}-\frac{36}{x}=1.\]
При $x>6$ оба знаменателя отличны от нуля. Приведём дроби к общему знаменателю:
\[\frac{36x-36(x-6)}{x(x-6)}=1,
\qquad \frac{216}{x(x-6)}=1.\]
Умножим обе части на ненулевой знаменатель:
\[x(x-6)=216\quad\Longleftrightarrow\quad x^2-6x-216=0.\]
Вычислим дискриминант и корни:
\[D=(-6)^2+4\cdot216=900,\qquad \sqrt D=30,\]
\[x_1=\frac{6+30}{2}=18,\qquad x_2=\frac{6-30}{2}=-12.\]
Отрицательный корень не удовлетворяет условию $x>6$: скорость автомобиля в этой задаче должна быть положительной.
\parttitle{Проверка} Скорость туда $18-6=12$ км/ч. Время туда $36/12=3$ ч, обратно $36/18=2$ ч. Разность составляет ровно 1 ч.
\result{Скорость на обратном пути равна 18 км/ч.}
\practice{Между пунктами 48 км. На обратном пути автомобиль едет на 4 км/ч быстрее и затрачивает на 1 ч меньше. Найдите скорость на обратном пути. Составьте таблицу и обоснуйте выбор корня.}
\vfill
\begin{center}\small\color{accent} Лёвин Артём Александрович эксклюзивно для Ярослава Верещагина\end{center}
\end{document}
